% LaTeX Template for Thesis Technical Report
% Optimized for academic submission

\documentclass[11pt,a4paper]{article}

% ============================================================================
% PACKAGES
% ============================================================================

% Typography and formatting
\usepackage[utf8]{inputenc}
\usepackage[T1]{fontenc}
\usepackage{geometry}
\geometry{
    a4paper,
    left=1.5in,
    right=1in,
    top=1in,
    bottom=1in
}

% Math and equations
\usepackage{amsmath}
\usepackage{amssymb}
\usepackage{amsfonts}
\usepackage{mathtools}
\usepackage{bm}  % Bold math symbols

% Graphics and figures
\usepackage{graphicx}
\usepackage{float}
\usepackage{subcaption}
\usepackage{tikz}
\usetikzlibrary{positioning,arrows.meta,shapes.geometric}

% Tables
\usepackage{booktabs}
\usepackage{multirow}
\usepackage{array}

% Code listings
\usepackage{listings}
\usepackage{xcolor}

% Algorithms
\usepackage{algorithm}
\usepackage{algorithmic}

% References and citations
\usepackage{hyperref}
\hypersetup{
    colorlinks=true,
    linkcolor=blue,
    filecolor=magenta,
    urlcolor=cyan,
    citecolor=green,
    pdftitle={Multi-Robot SLAM Mapping Module},
    pdfauthor={Piyush Bhansali}
}

% Header and footer
\usepackage{fancyhdr}
\pagestyle{fancy}
\fancyhf{}
\rhead{Multi-Robot SLAM Mapping Module}
\lhead{P. Bhansali}
\rfoot{Page \thepage}

% Section formatting
\usepackage{titlesec}
\titleformat{\section}{\Large\bfseries}{\thesection}{1em}{}
\titleformat{\subsection}{\large\bfseries}{\thesubsection}{1em}{}

% ============================================================================
% CODE LISTING STYLE
% ============================================================================

\definecolor{codegreen}{rgb}{0,0.6,0}
\definecolor{codegray}{rgb}{0.5,0.5,0.5}
\definecolor{codepurple}{rgb}{0.58,0,0.82}
\definecolor{backcolour}{rgb}{0.95,0.95,0.92}

\lstdefinestyle{pythonstyle}{
    backgroundcolor=\color{backcolour},
    commentstyle=\color{codegreen},
    keywordstyle=\color{magenta},
    numberstyle=\tiny\color{codegray},
    stringstyle=\color{codepurple},
    basicstyle=\ttfamily\footnotesize,
    breakatwhitespace=false,
    breaklines=true,
    captionpos=b,
    keepspaces=true,
    numbers=left,
    numbersep=5pt,
    showspaces=false,
    showstringspaces=false,
    showtabs=false,
    tabsize=2,
    language=Python
}

\lstset{style=pythonstyle}

% ============================================================================
% CUSTOM COMMANDS
% ============================================================================

% Math operators
\DeclareMathOperator{\argmin}{argmin}
\DeclareMathOperator{\argmax}{argmax}
\DeclareMathOperator{\diag}{diag}
\DeclareMathOperator{\trace}{tr}

% Special sets
\newcommand{\R}{\mathbb{R}}
\newcommand{\SE}{\text{SE}}

% Vectors and matrices
\newcommand{\vect}[1]{\boldsymbol{#1}}
\newcommand{\mat}[1]{\mathbf{#1}}

% Theorem environments
\usepackage{amsthm}
\newtheorem{theorem}{Theorem}[section]
\newtheorem{lemma}[theorem]{Lemma}
\newtheorem{proposition}[theorem]{Proposition}
\newtheorem{corollary}[theorem]{Corollary}
\theoremstyle{definition}
\newtheorem{definition}{Definition}[section]
\newtheorem{example}{Example}[section]
\theoremstyle{remark}
\newtheorem{remark}{Remark}

% ============================================================================
% DOCUMENT CONTENT
% ============================================================================

\begin{document}

% Title page
\begin{titlepage}
    \centering
    \vspace*{2cm}

    {\Huge\bfseries Multi-Robot SLAM Mapping Module\\[0.5cm]}
    {\Large Technical Report}

    \vspace{1.5cm}

    {\Large\itshape Piyush Bhansali}

    \vspace{1cm}

    {\large Thesis Work\\
    December 2025}

    \vfill

    {\large ROS2 Jazzy Multi-Robot SLAM Framework}

    \vspace{1cm}

    {\large \today}

\end{titlepage}

% Table of contents
\tableofcontents
\newpage

% Abstract
\begin{abstract}
This report presents a comprehensive multi-robot Simultaneous Localization and Mapping (SLAM) system designed for autonomous exploration in structured indoor environments. The system employs a submap-based mapping approach with hybrid feature extraction, multi-stage loop closure detection, and pose graph optimization using GTSAM. Key contributions include: (1) hybrid feature representation combining Scan Context with geometric features, (2) two-stage loop closure detection with coarse Scan Context filtering and fine geometric verification, (3) scale-invariant matching for robust place recognition, (4) selective loop closure with geometric distinctiveness filtering, and (5) real-time pose correction via Extended Kalman Filter integration. The system achieves sub-5cm positional accuracy with loop closure, 85\%+ loop closure detection rate on distinctive features, and operates in real-time at 10 Hz.
\end{abstract}

\newpage

% ============================================================================
% MAIN CONTENT GOES HERE
% Include the converted markdown content
% ============================================================================

$body$

% ============================================================================
% BIBLIOGRAPHY (if needed)
% ============================================================================

\bibliographystyle{IEEEtran}
% \bibliography{references}  % Uncomment if you have a .bib file

\end{document}
